% ============================================================================
% Chapter 13: Barrier-Limited Release Systematics
% ============================================================================
% Source: radioactivity_v7_1_alpha15/04_GN_FIT_AND_RESIDUALS.md
% Status: FILLED from SoT
% Contamination check: PASSED (no banned terms)
% Contamination guard: 3D-model terminology routed to Appendix X/Q.
% Layer A text uses EDC-native vocabulary only.
% ============================================================================

\chapter{Barrier-Limited Release Systematics}
\label{ch:barrier_release}

\source{radioactivity\_v7\_1\_alpha15}

% ----------------------------------------------------------------------------
\section*{Abstract}
% ----------------------------------------------------------------------------

This chapter establishes the \textbf{baseline scaling law} for barrier-limited
release of closed-4 units from heavy coordination clusters. The baseline lane
captures broad empirical regularities using two scalar parameters: a barrier
parameter and a release-energy parameter. We explicitly mark the baseline as
\textbf{[BL]}---it is not derived from EDC, but serves as the reference against
which EDC corrections are measured. The chapter identifies where the baseline
succeeds, where it systematically fails (high coordination frustration,
high-A region), and sets up the interface to Chapter~\ref{ch:frustrated} (coordination
frustration correction) and Chapter~\ref{ch:high_coord} (high-coordination predictions).

% ----------------------------------------------------------------------------
% CHAPTER SPINE
% ----------------------------------------------------------------------------
\paragraph{Chapter Spine.}
\textbf{Purpose:} Establish the baseline lane~[BL] for barrier-limited closed-4
release; define the residual interface for EDC corrections.
\textbf{Inputs:} (i)~Empirical release-time data~[BL],
(ii)~closed-4 as release unit from Ch.~\ref{ch:helium4}~[Der].
\textbf{Outputs:} (i)~Baseline law $\log t = aX + b$~[BL],
(ii)~residual definition $\Delta = \log(t_{\text{obs}}/t_{\text{BL}})$~[BL].
\textbf{Dependencies:} Ch.~\ref{ch:helium4} (closed-4 as fundamental release unit).
\textbf{Forward links:} Residuals to Ch.~\ref{ch:frustrated} (frustration correction);
predictions to Ch.~\ref{ch:high_coord}.
\textbf{Epistemic note:} This chapter is \emph{baseline only}~[BL]---not an EDC derivation.

% ============================================================================
\section{Motivation: A Baseline Law for Barrier-Limited Release}
\label{sec:gn:motivation}
% ============================================================================

\subsection{The Phenomenological Problem}

\begin{tcolorbox}[colback=gray!5,colframe=gray!50!black,title=Epistemic Tag: {[BL]} Baseline]
Throughout this chapter, the tag \tagBL{} marks \textbf{baseline} content:
empirical regularities or phenomenological fits that are \emph{not} derived
from EDC first principles. Baseline data serves as the reference against
which EDC predictions are compared. The baseline lane is descriptive, not
explanatory---it captures patterns without explaining their origin.
\end{tcolorbox}

Across a wide range of heavy coordination clusters (cluster size $A > 200$),
the release of a closed-4 unit follows a striking empirical regularity:
the logarithm of the release time correlates strongly with a simple combination
of two parameters:

\begin{enumerate}
    \item A \textbf{barrier parameter} related to the cluster's charge boundary
    \item A \textbf{release-energy parameter} related to the energy available
          for the transition
\end{enumerate}

This regularity spans more than 20 orders of magnitude in release times,
from sub-microsecond to cosmological timescales. \tagBL

\subsection{Purpose of This Chapter}

\begin{resultbox}
    \textbf{Scope declaration:} This chapter establishes the \emph{baseline lane}
    used for residual analysis. It is \textbf{not} an EDC derivation---the baseline
    formula is an empirical regularity tagged [BL]. EDC enters only through
    the correction term developed in Chapter~\ref{ch:frustrated}.
    \tagBL
\end{resultbox}

The baseline lane serves as the ``null hypothesis'' against which structural
effects (coordination frustration, boundary mismatch) are tested.

% ============================================================================
\section{Baseline Scaling Form}
\label{sec:gn:baseline}
% ============================================================================

\subsection{The Baseline Lane [BL]}

The empirical baseline law relates the release time $t_{1/2}$ to the barrier
and release-energy parameters:

\begin{derivationbox}
    \textbf{Baseline lane formula:}
    \begin{equation}
        \log_{10}(t_{1/2}) = a \times X + b
        \label{eq:baseline_lane}
    \end{equation}
    where the barrier-to-release ratio is:
    \begin{equation}
        X = \frac{Z_{\text{eff}}}{\sqrt{E_{\text{rel}}}}
        \label{eq:barrier_ratio}
    \end{equation}
    \tagBL
\end{derivationbox}

\subsection{EDC-Native Variable Definitions}

Table~\ref{tab:variables} defines the baseline variables in Book IV terminology.

\begin{table}[h]
\centering
\caption{Baseline variables dictionary (EDC-native).}
\label{tab:variables}
\begin{tabular}{lp{7cm}c}
    \toprule
    Symbol & EDC-Native Definition & Tag \\
    \midrule
    $t_{1/2}$ & Release time: characteristic timescale for closed-4
               emission from parent cluster & [BL] \\
    $Z_{\text{eff}}$ & Barrier parameter: effective charge determining
                       the height of the coordination boundary through
                       which the closed-4 unit tunnels & [BL] \\
    $E_{\text{rel}}$ & Release energy: kinetic energy available to the
                       emitted closed-4 unit after barrier penetration & [BL] \\
    $X$ & Barrier-to-release ratio: dimensionless combination
          $Z_{\text{eff}}/\sqrt{E_{\text{rel}}}$ & [BL] \\
    $a, b$ & Baseline coefficients (Appendix~Q) & [BL] \\
    \bottomrule
\end{tabular}
\end{table}

\subsection{Baseline Coefficients}

The coefficients $a$ and $b$ are determined by regression on empirical data.

\begin{quarantinebox}[title=Appendix Q: Coefficient Estimation]
    The fitting procedure, dataset selection, regression method, and
    coefficient values are documented in Appendix~Q. Layer A uses only
    the functional form; numerical calibration is quarantined.
\end{quarantinebox}

For reference, typical values are $a \approx 1.5$ and $b \approx -47$,
giving $R^2 \approx 0.99$ across representative datasets. \tagBL

\begin{quarantinebox}[title=Appendix X: Conventional Notation]
    In conventional literature, this scaling law is known by historical
    names that we do not use in Layer A. The mapping to conventional
    notation is provided in Appendix~X for reader orientation.
\end{quarantinebox}

% ============================================================================
\section{Physical Picture (EDC-Native)}
\label{sec:gn:picture}
% ============================================================================

\subsection{The Release Mechanism}

In EDC terms, barrier-limited release occurs when:

\begin{enumerate}
    \item A \textbf{closed-4 unit} (the minimal closed topological unit from
          Chapter~\ref{ch:helium4}) exists within a larger parent cluster
    \item The closed-4 unit tunnels through an \textbf{effective barrier}
          shaped by the parent cluster's coordination boundary
    \item The tunneling rate is controlled by two parameters:
          \begin{itemize}
              \item Barrier height $\propto Z_{\text{eff}}$
              \item Available energy $\propto E_{\text{rel}}$
          \end{itemize}
\end{enumerate}

\subsection{What the Baseline Ignores}

The baseline lane treats the parent cluster as \textbf{featureless} except
for two scalar parameters. It ignores:

\begin{itemize}
    \item Internal coordination structure (how junctions are arranged)
    \item Coordination frustration (mismatch between allowed and actual
          coordination numbers)
    \item Structural reconfiguration costs during emission
    \item Preferred vs.\ forbidden emission channels
\end{itemize}

These omitted effects generate the \textbf{residuals} that motivate
Chapter~\ref{ch:frustrated}. \tagP

\subsection{Why a 1D Effective Barrier Works}

The reduction to a single barrier-to-release ratio $X$ works because:

\begin{enumerate}
    \item The closed-4 unit is compact and rigid (Chapter~\ref{ch:helium4})
    \item The tunneling probability is exponentially sensitive to $X$
    \item Most structural details are ``averaged out'' by the exponential
\end{enumerate}

This 1D reduction is an approximation that captures $\sim 99\%$ of the
variance in $\log_{10}(t_{1/2})$---but the remaining $\sim 1\%$ contains
the coordination information. \tagP + \tagBL

% ============================================================================
\section{Why It Works (and Where It Must Fail)}
\label{sec:gn:failure}
% ============================================================================

\subsection{Success of the Baseline}

The baseline lane succeeds broadly because:

\begin{enumerate}
    \item \textbf{Exponential sensitivity:} Small changes in $X$ produce
          large changes in $t_{1/2}$, dominating over structural corrections
    \item \textbf{Universal emitter:} The closed-4 unit has fixed properties
          (Chapter~\ref{ch:helium4}), so emission physics is reproducible
    \item \textbf{Smooth barrier:} For most clusters, the coordination
          boundary is sufficiently regular that a single effective height
          captures the tunneling
\end{enumerate}

\begin{resultbox}
    \textbf{Baseline success:} The two-parameter lane explains $\gtrsim 98\%$
    of the variance in $\log_{10}(t_{1/2})$ across $\sim 20$ orders of magnitude.
    \tagBL
\end{resultbox}

\subsection{Failure Modes (EDC-Native)}

The baseline systematically fails when:

\begin{enumerate}
    \item \textbf{High coordination frustration:}
          When the parent cluster's coordination number $n(A)$ deviates
          significantly from the allowed set (Chapter~\ref{ch:frustrated}),
          structural strain modifies the effective barrier.
          \tagP

    \item \textbf{Boundary mismatch:}
          When the emitting closed-4 unit must cross a coordination boundary
          with mismatched local structure, additional reconfiguration energy
          is required.
          \tagP

    \item \textbf{Structural reconfiguration costs:}
          When emission requires the parent cluster to rearrange its
          junction network (not captured by baseline scalars).
          \tagP

    \item \textbf{Forbidden-zone cases:}
          When the parent cluster lies in a topologically ``forbidden''
          region where certain coordination numbers are inaccessible,
          the emission dynamics differ qualitatively.
          \tagP
\end{enumerate}

\subsection{Bridge to Chapter~\ref{ch:frustrated}}

\begin{resultbox}
    \textbf{Motivation for correction term:}
    Residuals from the baseline lane correlate with coordination frustration
    metrics. This motivates adding a structure-dependent correction in
    Chapter~\ref{ch:frustrated}:
    \begin{equation}
        \log_{10}(t_{1/2}) = a X + b + g \times \dfrustration(n)
        \label{eq:corrected_lane}
    \end{equation}
    where $\dfrustration(n)$ is the coordination frustration metric and $g$
    is a coupling coefficient.
    \tagP
\end{resultbox}

% ============================================================================
\section{Baseline Residual Definition}
\label{sec:gn:residual}
% ============================================================================

\subsection{Definition}

The \textbf{baseline residual} measures the deviation of observed release
times from baseline predictions:

\begin{derivationbox}
    \begin{equation}
        \Delta = \log_{10}(t_{\text{obs}}) - \log_{10}(t_{\text{BL}})
        \label{eq:residual_def}
    \end{equation}
    where $t_{\text{BL}}$ is the baseline prediction from Eq.~\eqref{eq:baseline_lane}.
    \tagBL
\end{derivationbox}

\subsection{Interpretation}

\begin{center}
\begin{tabular}{cp{10cm}}
    \toprule
    $\Delta$ & Interpretation \\
    \midrule
    $\Delta > 0$ & Observed release is \textbf{slower} than baseline
                   (``hindered'' emission) \\
    $\Delta < 0$ & Observed release is \textbf{faster} than baseline
                   (``favored'' emission) \\
    $\Delta \approx 0$ & Baseline adequately describes the system \\
    \bottomrule
\end{tabular}
\end{center}

\subsection{Interface to Chapter~\ref{ch:frustrated}}

Chapter~\ref{ch:frustrated} will show that:
\begin{equation}
    \Delta \approx g \times \dfrustration(n)
    \label{eq:delta_d}
\end{equation}

The coupling $g$ and the frustration metric $\dfrustration(n)$ are defined
and calibrated in Chapter~\ref{ch:frustrated}; any regression details are
routed to Appendix~Q. \tagP

% ============================================================================
\section{Tables}
\label{sec:gn:tables}
% ============================================================================

\subsection{Table 13.2: Baseline Formula and Epistemic Tags}

\begin{table}[h]
\centering
\caption{Baseline formula components with epistemic tags.}
\label{tab:formula}
\begin{tabular}{lcc}
    \toprule
    Component & Formula & Tag \\
    \midrule
    Baseline lane & $\log_{10}(t_{1/2}) = a X + b$ & [BL] \\
    Barrier ratio & $X = Z_{\text{eff}}/\sqrt{E_{\text{rel}}}$ & [BL] \\
    Coefficients & $a \approx 1.5$, $b \approx -47$ & [BL] (Appendix Q) \\
    Residual & $\Delta = \log_{10}(t_{\text{obs}}/t_{\text{BL}})$ & [BL] \\
    Correction (Ch.~\ref{ch:frustrated}) & $+ g \times \dfrustration(n)$ & [P] \\
    \bottomrule
\end{tabular}
\end{table}

\subsection{Table 13.3: Failure/Residual Taxonomy}

\begin{table}[h]
\centering
\caption{Baseline failure modes and EDC corrections.}
\label{tab:failure}
\begin{tabular}{p{4cm}p{5cm}cc}
    \toprule
    Failure Mode & EDC-Native Description & Chapter & Tag \\
    \midrule
    High coordination frustration &
    $n(A)$ deviates from allowed coordination set &
    \ref{ch:frustrated} & [P] \\
    Boundary mismatch &
    Coordination boundary incompatible with emission &
    \ref{ch:frustrated} & [P] \\
    Structural reconfiguration &
    Parent must rearrange junctions during emission &
    \ref{ch:frustrated} & [P] \\
    Forbidden-zone emission &
    Parent lies in topologically forbidden region &
    \ref{ch:high_coord} & [P] \\
    High-A anomalies &
    Extreme coordination frustration in heavy clusters &
    \ref{ch:high_coord} & [P] \\
    \bottomrule
\end{tabular}
\end{table}

% ============================================================================
\section{Epistemic Status and Falsifiability}
\label{sec:gn:epistemic}
% ============================================================================

\subsection{Epistemic Status Table}

\begin{table}[h]
\centering
\caption{Epistemic status of Chapter~\ref{ch:barrier_release} claims.}
\label{tab:gn:epistemic}
\begin{tabular}{p{5.5cm}cc}
    \toprule
    Claim & Tag & Confidence \\
    \midrule
    Baseline lane formula (Eq.~\ref{eq:baseline_lane}) & [BL] & Established \\
    Barrier ratio definition (Eq.~\ref{eq:barrier_ratio}) & [BL] & Established \\
    $R^2 \gtrsim 0.98$ for baseline fit & [BL] & Established \\
    1D reduction approximation & [P] & High \\
    Residuals encode coordination information & [P] & Medium--High \\
    Correction form $g \times \dfrustration$ & [P] & Medium \\
    Failure modes taxonomy & [P] & Medium \\
    \bottomrule
\end{tabular}
\end{table}

\subsection{Falsifiability Hooks}

\begin{edcopenbox}[title=Falsifiability Hooks]
    \begin{enumerate}
        \item \textbf{Residual-frustration correlation:}
              If baseline residuals $\Delta$ do \emph{not} correlate with
              the coordination frustration metric $\dfrustration(n)$,
              the EDC correction hypothesis fails.

        \item \textbf{Forbidden-zone outliers:}
              Clusters in topologically ``forbidden'' zones must show
              systematically large residuals under the baseline. If they
              follow the baseline lane, the forbidden-zone concept is wrong.

        \item \textbf{High-coordination systematic sign:}
              In the high-A region, residuals must show a consistent
              sign (faster or slower than baseline). Random scatter would
              invalidate the coordination frustration mechanism.

        \item \textbf{Closed-4 emission bias:}
              The baseline assumes emission of the \emph{minimal} closed unit.
              If alternative emission channels (closed-8, fragmentation)
              dominate, the baseline formulation breaks.
              (Forward-link to Chapter~\ref{ch:light_clusters} minimal-unit
              concept.)

        \item \textbf{Coefficient universality:}
              If baseline coefficients $a$, $b$ vary systematically across
              regions in a way not explained by EDC corrections, the
              single-lane picture fails.
    \end{enumerate}
\end{edcopenbox}

% ============================================================================
% Observer-side projection
% ============================================================================

\begin{observerbox}
Observer-side projection note: quoted terms are measurement labels for the
3D/4D projection of the 5D objects defined in this chapter; they do not
imply any conventional mechanism.

\begin{itemize}
    \item \ClosedFourRelease{} time $\leftrightarrow$ measured time proxy (projection label)
    \item Barrier-to-release ratio $\leftrightarrow$ empirical scaling variable (projection label)
\end{itemize}

What the observer records: the baseline lane $\log_{10}(t) = a \times X + b$
correlates 5D barrier parameters with measured release times. The residuals
$\Delta$ provide the interface to coordination corrections.
\end{observerbox}

% ============================================================================
\section{Summary and Forward Links}
\label{sec:gn:summary}
% ============================================================================

\begin{resultbox}
    \textbf{What this chapter provides:}
    \begin{itemize}
        \item Baseline lane for barrier-limited closed-4 emission [BL]
        \item EDC-native variable definitions (barrier parameter, release
              energy, barrier-to-release ratio)
        \item Physical picture: 1D effective barrier with two-parameter
              control
        \item Residual definition as interface to coordination corrections
        \item Failure mode taxonomy pointing to EDC structure terms
    \end{itemize}

    \textbf{What this chapter omits:}
    \begin{itemize}
        \item Derivation from EDC (baseline is empirical [BL])
        \item Numerical coefficient values (Appendix Q)
        \item Conventional terminology mapping (Appendix X)
        \item Coordination frustration correction (Chapter~\ref{ch:frustrated})
    \end{itemize}
\end{resultbox}

\bigskip
\noindent
\textbf{Forward references:}
\begin{itemize}
    \item Chapter~\ref{ch:frustrated}: Coordination frustration correction
          $g \times \dfrustration(n)$
    \item Chapter~\ref{ch:high_coord}: High-coordination predictions and extreme
          frustration
    \item Appendix~Q: Baseline regression details and coefficient calibration
    \item Appendix~X: Mapping to conventional notation
\end{itemize}

\bigskip
\noindent
\textbf{Derivation chain position:}
\begin{equation*}
    \text{Baseline lane [BL]}
    \xrightarrow{+ g \times \dfrustration}
    \text{Corrected lane (Ch.~\ref{ch:frustrated})}
    \xrightarrow{\text{predictions}}
    \text{High-A (Ch.~\ref{ch:high_coord})}
\end{equation*}

% ----------------------------------------------------------------------------
% BRIDGE TO NEXT CHAPTER
% ----------------------------------------------------------------------------
\paragraph{Bridge to Chapter~\ref{ch:frustrated}.}
The baseline law captures $> 98\%$ of variance in release times, but
systematic residuals remain. These residuals correlate with \emph{structure}:
clusters with ``awkward'' coordination numbers release faster than the
baseline predicts. Chapter~\ref{ch:frustrated} develops the coordination
frustration correction $\Delta = g \times d(n)$ using the $\Msix$ allowed
set from Chapter~\ref{ch:M6}.

